\documentclass[11pt]{article}
\usepackage[a4paper,margin=27mm]{geometry}
\usepackage{amsmath,amssymb,amsthm,booktabs,array}
\usepackage[T1]{fontenc}
\usepackage[colorlinks=true,linkcolor=blue,citecolor=blue,urlcolor=blue]{hyperref}
\hypersetup{pdftitle={Schatten and rank extensions of Clarisse and Wocjan's permutation classification},pdfauthor={Haobo Ma and Wenlin Zhang}}
\newtheorem{theorem}{Theorem}
\newtheorem{lemma}[theorem]{Lemma}
\newtheorem{corollary}[theorem]{Corollary}
\newcommand{\tr}{\operatorname{Tr}}
\newcommand{\ket}[1]{\lvert #1\rangle}
\newcommand{\bra}[1]{\langle #1\rvert}
\newcommand{\norm}[1]{\lVert #1\rVert}
\newcommand{\bits}{\{0,1\}}
\newcommand{\R}{\mathcal R}
\title{Schatten and rank extensions of Clarisse and Wocjan's\\permutation classification}
\author{Haobo Ma\\\small ChronoAI Pte Ltd, Singapore
\and Wenlin Zhang\\\small National University of Singapore\\\small ChronoAI Pte Ltd, Singapore}
\date{16 November 2026}
\begin{document}
\maketitle
\begin{abstract}
Extending Clarisse and Wocjan's orbit bound and answering their
norm-dependence question \cite[Theorem~3, p.~3, equation~(2);
Section~V, p.~6; Definition~1, p.~2]{CW},
arXiv:quant-ph/0504160v2 (14 June 2005), we prove for every $r\geq1$
and $d\geq2$ that mixed single-party and Bell-pair states distinguish
tensor-index permutations exactly up to the four row-set symmetries.
Consequently the prior orbit bound is exact for norm equality.
The same classification holds for rank and every Schatten functional
with $0<p\leq\infty$, $p\ne2$, including quasi-norms below one;
all permutations preserve the Hilbert--Schmidt norm.
We also give an adaptive identifying family of at most $2r-1$ states
and a Bell-state proof of the known partial-transpose incomparability.
\end{abstract}

\section{Statement and relation to the source}

Permutation separability criteria compare the trace norm of a rearranged
density matrix with one \cite{HHH}. Wocjan and Horodecki classified
dependencies obtained from norm-preserving maps on all operators
\cite[Lemmas 1--2]{WH}. Clarisse and Wocjan used Hermiticity to identify
further pairs, and counted the resulting orbits
\cite[Theorem 3]{CW}. Their remaining norm-dependence question concerns
whether positivity forces any additional identifications. We answer it
by an explicit family of positive trace-one separating states.
The earlier orbit count and partial-transpose independence are credited
below as prior results, rather than consequences first posed here.
The equality problem also appears as a probable class count in
OEIS A045723 \cite{OEIS}. Later work on random-state realignment
\cite[Section~1.1 and Theorem~1.1]{AN} and permutation moments
\cite{Liu} concerns detection and moment inequalities, respectively;
those results have a different scope from equality of norm functions.

Let $r\geq1$, $d\geq2$, and $H=(\mathbb C^d)^{\otimes r}$.
Number the tensor-index positions by
$P=\{0,\ldots,2r-1\}$: positions $2k$ and $2k+1$ are the ket and bra
indices of party $k\in\{0,\ldots,r-1\}$. For $\sigma\in S_{2r}$ use
the convention of \cite[Section II]{CW}, shifted to zero-based indices:
\[
 (\Lambda_\sigma\rho)_{i_0i_1,\ldots,i_{2r-2}i_{2r-1}}
 =\rho_{i_{\sigma(0)}i_{\sigma(1)},\ldots,
              i_{\sigma(2r-2)}i_{\sigma(2r-1)}}.
\]
All indices take values in $\{0,\ldots,d-1\}$.
Set $S_\sigma=\{p\in P:\sigma(p)\text{ is even}\}$, an $r$-element
subset of $P$. Let $C(S)=P\setminus S$ and let $J$ exchange $2k$ with
$2k+1$ at every party; $J(S)=\{J(p):p\in S\}$.
For a matrix $A$, $\norm{A}_1=\tr\sqrt{A^*A}$ is its trace norm.
Norm dependence means equality of $\norm{\Lambda_\sigma\rho}_1$ and
$\norm{\Lambda_\mu\rho}_1$ for every $\rho\geq0$ with $\tr\rho=1$.

\begin{theorem}\label{thm:main}
For $r\geq1$ and $d\geq2$, permutations $\sigma,\mu\in S_{2r}$ are
norm dependent if and only if
\begin{equation}\label{eq:orbit}
 S_\mu\in\{S_\sigma,C(S_\sigma),J(S_\sigma),CJ(S_\sigma)\}.
\end{equation}
\end{theorem}

The theorem concerns equality of norm functions. The source also considers
the stronger property that no entanglement-detection set is contained in
another. That detection-set comparison is a separate assertion.

\section{A finite family of separating states}

For $S\subset P$ with $|S|=r$, let $\R_S\rho$ be the matrix obtained by
placing the original indices in $S$ on the row side, and those in $C(S)$
on the column side, each in increasing order. Then $\Lambda_\sigma\rho$
and $\R_{S_\sigma}\rho$ differ by independent row and column permutations,
so have the same singular values. Define the labels
\[
 x_k(S)=(a_k,b_k)
   =(\mathbf1_{2k\in S},\mathbf1_{2k+1\in S})\in\bits^2,
 \qquad
 D(x,y)=\mathbf1_{x_1\ne y_1}+\mathbf1_{x_2\ne y_2}.
\]
Here a label bit is one precisely when its original position becomes a
row index. Call the labels equal-side if $a_k=b_k$.

Let $P_0=\ket0\bra0$. At party $k$ put the maximally mixed qubit
$Q=(\ket0\bra0+\ket1\bra1)/2$, and put $P_0$ at the other parties;
denote the product state by $M_k$. For distinct $k,l$ put
$\ket{\Phi}_{kl}=(\ket{00}+\ket{11})/\sqrt2$ at those two parties,
with $P_0$ elsewhere, and denote its density matrix by $B_{kl}$.
These are positive trace-one states in every $d\geq2$.

\begin{lemma}\label{lem:norms}
The nonzero singular values of their rearrangements are as follows:
\begin{center}
\begin{tabular}{lll}
\toprule
State & Label condition & Nonzero singular values\\
\midrule
$M_k$ & $a_k=b_k$ & $1/\sqrt2$\\
$M_k$ & $a_k\ne b_k$ & $1/2,\ 1/2$\\
$B_{kl}$ & $D(x_k,x_l)=0$ & $1$\\
$B_{kl}$ & $D(x_k,x_l)=1$ & $1/\sqrt2,\ 1/\sqrt2$\\
$B_{kl}$ & $D(x_k,x_l)=2$ & $1/2,\ 1/2,\ 1/2,\ 1/2$\\
\bottomrule
\end{tabular}
\end{center}
In particular, the mixed-state trace norms distinguish equal-side from
opposite-side labels, and the Bell-state trace norm is $2^{D(x_k,x_l)/2}$.
\end{lemma}

\begin{proof}
The two nonzero entries of $M_k$ have coefficient $1/2$. If its ket and
bra positions are on the same side, these entries form one row or column,
of Euclidean length $1/\sqrt2$. Otherwise they occupy distinct rows and
distinct columns and give two singular values $1/2$.

For $B_{kl}$ the four entries have coefficient $1/2$, with local ket
indices $(u,u)$ and local bra indices $(v,v)$ for $u,v\in\bits$.
The placement of these two bits proves every row of the table.
If the two labels agree and are equal-side, all four entries form one
row or column. If they agree and are opposite-side, they form a
$2\times2$ all-ones block divided by two. Both cases have one singular
value one. At distance two, each bit occurs on each side and the four
entries have distinct rows and distinct columns; the block is a
$4\times4$ permutation matrix divided by two. At distance one, one label
is equal-side and the other is opposite-side. One side then carries a
single bit, while the other carries both bits. The block, or its
transpose, consists of two orthogonal rows with two entries $1/2$ in
each. Its two singular values are $1/\sqrt2$. Unused coordinates add
zero rows and columns, so none of these values depends on $r$ or $d$.
\end{proof}

\begin{lemma}\label{lem:labels}
Two families $x_k,y_k\in\bits^2$ with the same equal-side parties and
the same distances $D(x_k,x_l)=D(y_k,y_l)$ for all $k,l$ differ by one
global coordinate exchange and one global simultaneous bit complement.
\end{lemma}

\begin{proof}
Split the parties into the equal-side sector and its complement.
Within either sector a label is determined by its first bit, and the
distance of two labels is zero or two according as those first bits
agree or disagree. Thus the equality relation on first bits agrees for
the two families in each sector. In a nonempty sector choose an anchor.
The change in its first bit determines a bit $c$; agreement with the
anchor shows that every first bit in that sector changes by $c$.
For an empty sector choose $c=0$.

Write $c_E,c_O$ for the two sector changes. Simultaneous complement
changes both sectors' first bits, whereas coordinate exchange changes
the first bit only in the opposite-side sector. Consequently coordinate
exchange $c_E\mathbin\oplus c_O$ times followed by complement $c_E$
times maps every $x_k$ to $y_k$. The sector fixes each second bit from
its first, proving the assertion, including empty sectors.
\end{proof}

\begin{proof}[Proof of Theorem~\ref{thm:main}]
Suppose the norm functions agree. Testing $M_k$ and using
Lemma~\ref{lem:norms} identifies the equal-side parties. Testing $B_{kl}$
identifies all pair distances since $1,\sqrt2,2$ are distinct.
Lemma~\ref{lem:labels} now gives \eqref{eq:orbit}: coordinate exchange
is $J$, and bit complement is $C$.

Conversely, equal row sets give independent row and column reindexing.
Complementing a row set transposes the rearranged matrix, up to such
reindexing. For Hermitian $\rho$, exchanging all original ket and bra
positions replaces every tensor entry by its complex conjugate. Thus
$\R_{J(S)}\rho$ is an entrywise conjugate of $\R_S\rho$, up to row and
column reindexing. Transpose, entrywise conjugation, and row and column
permutations preserve singular values. Each of the four sets in
\eqref{eq:orbit} therefore has the same trace norm on every state.
\end{proof}

\begin{corollary}\label{cor:finite}
For fixed $r\geq1$, $d\geq2$, equality of two permutation norm functions
on the $r+\binom r2$ states $M_k,B_{kl}$ is equivalent to equality on
every density matrix. The same four symmetries also characterize
equality of the full singular-value lists on every density matrix.
\end{corollary}
\begin{proof}
The necessity argument uses exactly those states. The converse argument
preserves each singular value, not just their sum.
\end{proof}

\begin{corollary}[Adaptive identification]\label{cor:adaptive}
For two given permutations it suffices to compare $r$ mixed-state
values and at most $r-1$ Bell-state values. If both equal-side sectors
are nonempty, at most $2r-2$ values in total suffice.
\end{corollary}
\begin{proof}
First compare every $M_k$. If a value differs, the functions differ.
Otherwise the common equal-side partition is known. In each nonempty
sector choose an anchor and compare its Bell value with every other
party of that sector. These distances recover all first bits relative
to the anchor, exactly as in Lemma~\ref{lem:labels}. Distances across
sectors are always one. With $s\in\{1,2\}$ nonempty sectors the number
of Bell comparisons is $r-s$, giving the stated bounds. The choices
depend on the mixed-state answers; this is not a fixed universal
family of $2r-1$ states.
\end{proof}

\section{Counting the classes and changing the norm}

\begin{corollary}\label{cor:count}
The exact number of norm-dependence classes, including the identity
criterion, is
\begin{equation}\label{eq:count}
 N_r=\frac14\left(\binom{2r}{r}+2^r+
       \mathbf1_{2\mid r}\binom r{r/2}\right),\qquad r\geq1.
\end{equation}
There are $N_r-1$ nontrivial classes. For $r=1,2,3,4$ these numbers
are respectively $0,2,6,22$.
\end{corollary}
\begin{proof}
Every $r$-subset is $S_\sigma$ for some permutation. Apply the orbit
counting lemma to the four-element group generated by $C,J$ on these
subsets. The identity fixes $\binom{2r}{r}$ subsets. Complement fixes
none, since $P$ is nonempty. A $J$-fixed subset is a union of whole
ket--bra pairs: it exists only for even $r$, and then one chooses $r/2$
of the $r$ pairs. Finally $CJ(S)=S$ means that exactly one position of
each pair lies in $S$, giving $2^r$ choices. Averaging these fixed-point
counts proves \eqref{eq:count}. The orbit of the even positions is
the identity class, whose norm is one on every state, so subtracting it
leaves $N_r-1$ classes.
\end{proof}

The formula is the upper bound already obtained in
\cite[Theorem 3, equation (2)]{CW}; the separating states make it exact
for norm dependence. In particular the small-party counts agree with
those in \cite[Section IV]{CW} without extending their detection-set
claims to other dimensions or more parties.

If $s_j(A)$ are the singular values of $A$, for $0<p<\infty$ write
$\norm A_p=(\sum_j s_j(A)^p)^{1/p}$, and let
$\norm A_\infty=\max_j s_j(A)$. This is a quasi-norm when $p<1$.
\begin{corollary}\label{cor:schatten}
Theorem~\ref{thm:main} and Corollary~\ref{cor:finite} hold with the
trace norm replaced by any Schatten functional, $0<p\leq\infty$,
$p\ne2$. Corollary~\ref{cor:adaptive} holds in the same range.
For $p=2$ all tensor-index permutations have the same norm
on every operator.
\end{corollary}
\begin{proof}
Put $1/p=0$ when $p=\infty$. Lemma~\ref{lem:norms} gives
\[
 \norm{\R_S M_k}_p=
 \begin{cases}2^{-1/2}&a_k=b_k,\\2^{1/p-1}&a_k\ne b_k,\end{cases}
 \qquad
 \norm{\R_S B_{kl}}_p=2^{D(x_k,x_l)(1/p-1/2)}.
\]
For $p\ne2$ the two mixed-state values are distinct, as are the three
Bell-state values. The same label reconstruction applies, and the
four symmetries preserve every singular value; the argument uses no
triangle inequality, so it also applies below one. At $p=2$ the squared norm
is the sum of squared absolute values of all entries, which an index
permutation leaves unchanged.
\end{proof}

Permutation moments also appear in experimentally accessible
entanglement criteria \cite{Liu}. Corollary~\ref{cor:schatten}, for
example at $p=4$, classifies equality of those fourth singular-value
moments as functions of the state; it is not a comparison of the
detection powers of moment-based inequalities.

\begin{corollary}[Rank and positive-definite tests]\label{cor:rank}
Two permutations have equal ranks on every density matrix if and only
if \eqref{eq:orbit} holds. For every fixed $0<p\leq\infty$, $p\ne2$,
equality of Schatten functionals on positive-definite density matrices
already implies equality on all density matrices.
\end{corollary}
\begin{proof}
The ranks on $M_k$ are respectively one and two in the equal-side
and opposite-side cases. On $B_{kl}$ the rank is $2^{D(x_k,x_l)}$.
They therefore recover the same label data, and the four symmetries
preserve rank. For the second assertion, approximate any density matrix
$\rho$ by $(1-\varepsilon)\rho+\varepsilon I/d^r$, with
$0<\varepsilon<1$, and use continuity of each singular value and of
the Schatten functional as $\varepsilon\downarrow0$. This continuity
argument concerns Schatten functionals, not the discontinuous rank.
\end{proof}

\begin{corollary}[Uncorrelated spectators]\label{cor:spectators}
If an additional party in density matrix $\tau$ is left unchanged,
then for every $0<p\leq\infty$,
\[
 \norm{(\Lambda_\sigma\rho)\otimes\tau}_p
   =\norm{\Lambda_\sigma\rho}_p\norm{\tau}_p.
\]
Thus the trace norm is unchanged for every density matrix $\tau$;
every Schatten functional is unchanged when $\tau$ is pure.
\end{corollary}
\begin{proof}
The singular values of a tensor product are all pairwise products of
the two singular-value lists, by tensoring singular-value decompositions.
Sum their $p$th powers, or take their maximum when $p=\infty$.
For density matrices $\norm\tau_1=1$, and for pure ones $\norm\tau_p=1$.
The spectator is a separate, unrearranged tensor factor here.
\end{proof}
The trace-norm case explains the uncorrelated-ancilla construction in
\cite[Sections IV.B--C]{CW}; it does not assert invariance under every
possible regrouping of a mixed ancillary party.

\begin{corollary}[Partial-transpose criteria]\label{cor:cuts}
For $r\geq1$ and $d\geq2$ there are $2^{r-1}-1$ nontrivial
partial-transpose criteria, up to global transpose. Their
entanglement-detection sets are pairwise incomparable.
\end{corollary}
\begin{proof}
Let $T_A$ transpose exactly the parties of a nonempty proper subset
$A$. It is equivalent to $T_{A^c}$ by global transpose.
The $2^r-2$ nontrivial subsets thus give $2^{r-1}-1$ cuts.
On $B_{kl}$, $\norm{T_A B_{kl}}_1$ is two if exactly one of $k,l$
belongs to $A$, and one otherwise, by Lemma~\ref{lem:norms}.
For inequivalent nontrivial cuts $A,B$, some pair is separated by
$A$ but not $B$: otherwise each block of the $B$-cut would be contained
in a block of the $A$-cut, forcing the two nontrivial two-block
partitions to coincide. Reversing $A,B$ gives a pair separated only
by $B$. Their Bell states witness both directions of noncontainment.
The assertion is vacuous when fewer than two such cuts exist.
\end{proof}
This supplies an explicit proof of the loop-only independence discussed
by Clarisse and Wocjan \cite[Section~V, p.~6]{CW}. It does not establish
incomparability between criteria containing realignments. In particular,
their Section~IV.A distinguishes the two-qubit case from dimensions
where partial transpose and realignment have incomparable detection sets.

\section{Coverage and verification scope}

\begin{center}\small
\begin{tabular}{@{}p{.36\linewidth}p{.57\linewidth}@{}}
\toprule
Source item in \cite{CW} & Coverage\\
\midrule
Definition 1, p.~2; Section V, p.~6: further norm dependencies? &
Answered: no. The four symmetries are complete for all $r\geq1$,
$d\geq2$ (Theorem~\ref{thm:main}).\\[3pt]
Theorem 3: four symmetries and orbit bound &
Reproved; prior proof Clarisse--Wocjan (2005 preprint/2006 journal).
The upper bound becomes exact for norm equality
(Corollary~\ref{cor:count}).\\[3pt]
Section IV: $2,6,22$ nontrivial classes &
Reproved for norm equality; prior proof Clarisse--Wocjan (2005).
No new detection-set conclusion follows from these counts.\\[3pt]
Section V: independence of loop-only criteria &
Reproved; prior proof Clarisse--Wocjan (2005).
Corollary~\ref{cor:cuts} gives both Bell-state witnesses.\\[3pt]
Sections IV.B--C: uncorrelated ancillas &
Prior trace-norm invariance reproved; Schatten product formula
in Corollary~\ref{cor:spectators}.\\[3pt]
Section V: realignment vs partial-transpose detection &
Prior proof cited (Clarisse--Wocjan, 2005); not reproved here.\\[3pt]
Abstract; Sections IV--V: all-orbit detection-set incomparability &
Not covered in general; not implied by norm inequality.\\[3pt]
Equation (2); OEIS A045723, Clarisse's comment (2006) &
Probable norm-class count made exact; the orbit formula is prior work.\\[3pt]
Schatten quasi-norms, rank, adaptive tests, full-rank norm tests &
Not posed in the source; proved in
Corollaries~\ref{cor:adaptive}, \ref{cor:schatten}, \ref{cor:rank}.\\
\bottomrule
\end{tabular}
\end{center}

\begin{flushleft}\small
Declaration names below are relative to the namespace\par
\texttt{D5.S3.Quantum.Entanglement.PermutationCriteria}.\par
The Lean declaration \texttt{ClarisseWocjanClassification.result} in the
\href{https://github.com/the-omega-institute/trureturing/blob/3464bc4487b2236dab97aaf7bbed237f0db28b71/D5/S3/Quantum/Entanglement/PermutationCriteria/ClarisseWocjanClassification.lean}{pinned classification source}
states exactly the equivalence of Theorem~\ref{thm:main}, with positive
trace-one inputs and $r\geq1$, $d\geq2$.
At that same revision, the declarations
\href{https://github.com/the-omega-institute/trureturing/blob/3464bc4487b2236dab97aaf7bbed237f0db28b71/D5/S3/Quantum/Entanglement/PermutationCriteria/LabelRecovery.lean}{\texttt{LabelRecovery.recover\_labels}}
and
\href{https://github.com/the-omega-institute/trureturing/blob/3464bc4487b2236dab97aaf7bbed237f0db28b71/D5/S3/Quantum/Entanglement/PermutationCriteria/RowSetRecovery.lean}{\texttt{RowSetRecovery.recover\_rowSets}}
give Boolean-label and row-set reconstruction.
\href{https://github.com/the-omega-institute/trureturing/blob/3464bc4487b2236dab97aaf7bbed237f0db28b71/D5/S3/Quantum/Entanglement/PermutationCriteria/Witnesses.lean}{\texttt{dependent\_sameSide}}
recovers the sector data; the
\href{https://github.com/the-omega-institute/trureturing/blob/3464bc4487b2236dab97aaf7bbed237f0db28b71/D5/S3/Quantum/Entanglement/PermutationCriteria/BellBlocks.lean}{\texttt{BellBlocks.traceNorm\_fourEntries}}
gives the trace norm of four correlated half entries from their coordinate
kernels, used in
\href{https://github.com/the-omega-institute/trureturing/blob/3464bc4487b2236dab97aaf7bbed237f0db28b71/D5/S3/Quantum/Entanglement/PermutationCriteria/BellWitnesses.lean}{\texttt{dependent\_pairDistance}}
to recover pair distances.
The count, adaptive identification, Schatten and rank extensions,
positive-definite testing, spectator formula and cut witnesses have written proofs above;
they are not the proposition of the classification declaration.
The accompanying
\texttt{verify.py} checks the integer Gram matrices of the local witness
blocks and enumerates the identifying signatures and symmetry orbits
for $1\leq r\leq8$. It also checks adaptive signatures through eight
parties and cut-separation witnesses through eight parties.
These are finite checks of the formulas, rather than the all-parameter proof.
\end{flushleft}

\begin{thebibliography}{9}
\bibitem{CW}
L.~Clarisse and P.~Wocjan,
\emph{On independent permutation separability criteria},
Quantum Information \& Computation \textbf{6}(3), 277--288 (2006).
\href{https://doi.org/10.26421/QIC6.3-4}{doi:10.26421/QIC6.3-4}.
Posing version: \href{https://arxiv.org/abs/quant-ph/0504160v2}{arXiv:quant-ph/0504160v2} (2005).
\bibitem{WH}
P.~Wocjan and M.~Horodecki,
\emph{Characterization of combinatorially independent permutation
separability criteria}, Open Systems \& Information Dynamics
\textbf{12}(4), 331--345 (2005).
\href{https://doi.org/10.1007/s11080-005-4483-2}{doi:10.1007/s11080-005-4483-2};
\href{https://arxiv.org/abs/quant-ph/0503129v2}{arXiv:quant-ph/0503129v2}.
\bibitem{HHH}
M.~Horodecki, P.~Horodecki and R.~Horodecki,
\emph{Separability of mixed quantum states: linear contractions and
permutation criteria}, Open Systems \& Information Dynamics
\textbf{13}(1), 103--111 (2006).
\href{https://doi.org/10.1007/s11080-006-7271-8}{doi:10.1007/s11080-006-7271-8};
\href{https://arxiv.org/abs/quant-ph/0206008}{arXiv:quant-ph/0206008}.
\bibitem{Liu}
Z.~Liu, Y.~Tang, H.~Dai, P.~Liu, S.~Chen and X.~Ma,
\emph{Detecting entanglement in quantum many-body systems via permutation
moments}, Phys. Rev. Lett. \textbf{129}, 260501 (2022).
\href{https://doi.org/10.1103/PhysRevLett.129.260501}{doi:10.1103/PhysRevLett.129.260501};
\href{https://arxiv.org/abs/2203.08391v2}{arXiv:2203.08391v2}.
\bibitem{AN}
G.~Aubrun and I.~Nechita,
\emph{Realigning random states}, J. Math. Phys. \textbf{53}, 102210 (2012).
\href{https://doi.org/10.1063/1.4759115}{doi:10.1063/1.4759115};
\href{https://arxiv.org/abs/1203.3974v1}{arXiv:1203.3974v1}.
\bibitem{OEIS}
OEIS Foundation Inc., \emph{The On-Line Encyclopedia of Integer Sequences},
\href{https://oeis.org/A045723}{A045723}, comment by L.~Clarisse,
28 April 2006; accessed 16 November 2026.
\end{thebibliography}

\medskip\noindent\small
\textcopyright\ 2026 Haobo Ma and Wenlin Zhang.
Licensed under the \href{https://creativecommons.org/licenses/by/4.0/}{Creative Commons Attribution 4.0 International License (CC BY 4.0)}.
\end{document}
